\section{如何阅读后续 CS25：四个接口与一张检查表}

导论的最终价值不是把历史停在 GPT-3/BERT，而是提供后续课程的比较语言。无论讲者讨论视觉、RL、生成模型、biology 或更大语言模型，都可以用 token interface、interaction pattern、training objective 与 deployment contract 四个问题拆解。

\subsection{四接口阅读法}

\begin{importantbox}{每讲都回答四个问题}
\begin{enumerate}
  \item \textbf{Token interface}：输入被切成什么单位，position 与 modality 如何表示？
  \item \textbf{Interaction pattern}：full、local、causal、cross、sparse 还是 recurrent？
  \item \textbf{Training objective}：next-token、mask、contrastive、control、ranking 还是 multi-task？
  \item \textbf{Deployment contract}：输出用于 generation、representation、decision、retrieval 还是 action？如何评估失败？
\end{enumerate}
\end{importantbox}

\teachervoice{三位讲师最后把观众导向完整 guest-speaker 系列。课堂提示：导论只给共同 vocabulary，真正的学习发生在后续讲者如何修改 token、mask、objective 和 system boundary。}

\begin{knowledgebox}{研究问题模板}
先选一个领域失败：长上下文、低数据、实时延迟、因果决策或安全约束；再指出标准 Transformer 哪个接口造成限制；提出可验证改动；给出 baseline、ablation 与资源预算；最后说明图表不能证明什么。这样“Transformer 应用”才从换数据集升级为研究设计。
\end{knowledgebox}

\subsection{本章小结}

后续 CS25 不应被读成模型目录。四接口方法帮助读者区分通用骨架、领域适配、训练信号与产品责任，并把每场 guest talk 转成可比较的系统设计。

